\section{Active Inference Primitives}\label{sec:active_inference_background}

PCE's design language is borrowed from the active-inference / free-energy programme \citep{Friston2010FreeEnergyUnifiedBrainTheory,Friston2017ActiveInferenceCuriosity}, but its v0.4 implementation does not yet realise that programme in its strict variational form. This section states the formal primitives, then declares precisely what the v0.4 cascade implements and what it does not. The v0.4 prose departs from the v0.3 prose mainly in the latter: the v0.4 release writes the scope-cuts down explicitly, in line with the mechanism-study framing.

\subsection{The formal substrate}

Let $q(\theta)$ denote a posterior over a candidate-discriminating parameter $\theta$ and $p(\theta\mid m)$ a model with prior $p(\theta\mid m)$. Free-energy minimisation~\citep{Friston2017ActiveInferenceCuriosity,Friston2010FreeEnergyUnifiedBrainTheory} selects models / candidates that minimise
\begin{equation}
F(q,m) = \mathbb{E}_{q(\theta)}\!\left[\log q(\theta) - \log p(\theta, o\mid m)\right],
\end{equation}
where $o$ is the observed outcome (the surface). \emph{Bayesian model reduction} \citep{FristonPenny2011,Friston2018BayesianModelReductionArxiv,FristonEtAl2016BayesianModelReductionEmpiricalBayesDCM} reduces a rich model $m_F$ to a sparser model $m_R$ whose evidence is
\begin{equation}
\log p(o\mid m_R) = \log p(o\mid m_F) + \log \int q(\theta\mid o, m_F)\,\frac{p(\theta\mid m_R)}{p(\theta\mid m_F)} \,\mathrm{d}\theta.
\end{equation}
Practically, in PCE we use this in closed form for nested Dirichlet--multinomial models over candidate categories, computing $\Delta F = F(q\mid m_F) - F(q\mid m_R)$ via the standard log-Gamma identity. A negative $\Delta F$ indicates posterior switch is supported by the data; a positive $\Delta F$ supports the prior.

\subsection{The link to \iast{apohana}}

\iast{Apohana-\'sakti} (``contrastive exclusion'') is precisely what BMR's reduced posterior performs: it excludes the candidates that the data does not support. We therefore identify the \iast{j\~n\=ana} stage of the cascade with BMR posterior selection across the $K$ candidates emitted by \iast{icch\=a}, gated by the \iast{apohana} contrastive geometry computed in embedding space. This identification is structural and not metaphorical: the operator prunes a finite candidate set against an aspect-conditioned prior, and its scoring head emits a ranking that is a posterior in the sense of an ordered density over $K$ atoms.

\subsection{Insight as posterior re-organisation}

\citet{Friston2017ActiveInferenceCuriosity} formalise insight as the BMR-driven re-organisation of $q(\theta)$ that yields a sudden drop in $F$. This is the active-inference correlate of \iast{vimar\'sa}: the recursive layer fires when the post-\iast{kriy\=a} state's $F$ has not yet stabilised, indicating an unrecognised aspect of the constraint manifold. In v0.4 the \iast{vimar\'sa} operator (Appendix~\ref{app:operator_specs}) conditions on three signals jointly: aspect-multiplicity in the surface (cosine to multiple aspect targets), novelty against the storehouse (Hopfield retrieval similarity), and the BMR $\Delta F$. The result of H8b in this paper is that the disjunctive event-gate as currently defined under-fires; an H8b-informed v0.5 will move to a \emph{learned} gate trained on the H8a paired-revision data (already supported by \texttt{commit\_policy="learned\_gate"}, \S\ref{sec:methods}, ADR-002).

\subsection{What v0.4 implements honestly}

The v0.4 cascade implements an \emph{LM-best-of-K composite-score} substitute for true variational free-energy minimisation. The composite is a weighted sum of axis scorers (see \S\ref{sec:benchmarks}); the per-candidate best-of-K is taken over $K$ samples from the \iast{icch\=a} fan-out. We treat the \emph{rank order} this composite produces as a stand-in for $q(\theta)$, and the per-candidate score gap as a stand-in for the $\Delta F$ that BMR would compute. This is a design choice with three honest costs: (i) the composite is \emph{learned} only insofar as the weights were tuned on a small held-out sample, not by minimising any tractable variational lower bound; (ii) what we call $\Delta F$ in the v0.4 traces is therefore a score-gap, not a true free-energy delta; (iii) the per-item \texttt{FreeEnergyBudget} (\texttt{audit/v0.3/} ledger; ADR-003 v0.4) gates the cascade's token spend on a \emph{token-cost} budget that is interpreted as a free-energy budget at the cost-accounting level only. None of this invalidates the architecture --- best-of-K is a well-studied surrogate for variational selection in the LM literature \citep{StiennonEtAl2020LearningToSummarizeFromHumanFeedback,MadaanEtAl2023SelfRefineNeurIPS} --- but the prose has to say so. The v0.5 ladder commits to lifting the substitute to a true variational scoring head against a discrete-state-space model in \texttt{pymdp v1.0}~\citep{Heins2022pymdp,pymdpRelease100}.

\subsection{The Hopfield \iast{\=alayavij\~n\=ana}: status in v0.4}

The \iast{\=alayavij\~n\=ana} (``store-consciousness'') in PCE is implemented as a Hopfield-style attractor network~\citep{KrotovHopfield2016DenseAssociativeMemory,RamsauerEtAl2020HopfieldNetworksIsAllYouNeed,Mtenga2024HopfieldCreativity,Weber2025SelfOptHopfield} read by \iast{apohana} as a warm-start aspect prior and written by \iast{vimar\'sa.consolidate} after every committed surface (Appendix~\ref{app:consolidation}). The v0.4 honest claim is narrower than the v0.3 prose suggested: the Hopfield store is wired and exercised in single-session runs, but no large-scale multi-session memory dynamics have been measured against the network's documented retrieval guarantees \citep{RamsauerEtAl2020HopfieldNetworksIsAllYouNeed}. The v0.4 ablation stack treats the Hopfield store as one of three signals to the \iast{vimar\'sa} gate; we have not yet measured the marginal contribution of the Hopfield-warm-start signal alone. This is on the v0.5 ladder.
